\documentclass[11pt,a4paper]{article}

% ---------- Sprache & Layout ----------
\usepackage[utf8]{inputenc}
\usepackage[T1]{fontenc}
\usepackage[ngerman]{babel}
\usepackage[a4paper,margin=2.4cm]{geometry}
\usepackage{microtype}
\usepackage{parskip}
\usepackage{enumitem}
\usepackage{array}
\usepackage{booktabs}
\usepackage{xcolor}
\usepackage{hyperref}
\usepackage{listings}

% ---------- Hyperlinks ----------
\hypersetup{
  colorlinks=true,
  linkcolor=black,
  urlcolor=blue,
  pdftitle={Restrukturierte C++- und OOP-Notizen},
  pdfauthor={}
}

% ---------- Code-Layout ----------
\definecolor{codegray}{gray}{0.95}
\lstdefinestyle{cppstyle}{
  language=C++,
  backgroundcolor=\color{codegray},
  basicstyle=\ttfamily\small,
  keywordstyle=\bfseries,
  commentstyle=\itshape,
  stringstyle=\ttfamily,
  numbers=left,
  numberstyle=\tiny,
  stepnumber=1,
  numbersep=8pt,
  frame=single,
  breaklines=true,
  tabsize=2,
  showstringspaces=false,
  columns=fullflexible,
  keepspaces=true,
  inputencoding=utf8,
  extendedchars=true,
  literate={ä}{{\"a}}1 {ö}{{\"o}}1 {ü}{{\"u}}1 {Ä}{{\"A}}1 {Ö}{{\"O}}1 {Ü}{{\"U}}1 {ß}{{\ss}}1
}
\lstset{style=cppstyle}

\newcommand{\code}[1]{\texttt{#1}}
\newcommand{\important}[1]{\textbf{#1}}

\title{Restrukturierte Notizen zu C++ und objektorientierter Programmierung}
\author{}
\date{}

\begin{document}
\maketitle
\tableofcontents
\newpage

\section{Objektorientierte Programmierung}

\subsection{Die vier Grundsäulen der OOP}
\begin{enumerate}[label=\arabic*.]
  \item \important{Vererbung}: Klassen können Eigenschaften und Verhalten anderer Klassen übernehmen.
  \item \important{Abstraktion}: Wesentliche Eigenschaften werden hervorgehoben, unwichtige Details werden ausgeblendet.
  \item \important{Polymorphie}: Objekte können über gemeinsame Schnittstellen unterschiedlich reagieren.
  \item \important{Datenkapselung}: Daten werden geschützt und kontrolliert über Methoden zugänglich gemacht.
\end{enumerate}

\subsection{Polymorphie}
\subsubsection{Statische Polymorphie}
Statische Polymorphie wird zur Compilezeit aufgelöst. Typische Beispiele sind Funktionsüberladung und Templates.

\subsubsection{Dynamische Polymorphie}
Dynamische Polymorphie wird zur Laufzeit aufgelöst. Sie entsteht typischerweise durch:
\begin{itemize}
  \item Vererbung,
  \item virtuelle Funktionen,
  \item Aufrufe über Basisklassen-Zeiger oder Basisklassen-Referenzen.
\end{itemize}

\begin{lstlisting}[caption={Dynamische Polymorphie mit virtueller Funktion}]
class Base {
public:
    virtual void print() const;
    virtual ~Base() = default;
};

class Child : public Base {
public:
    void print() const override;
};

void call(const Base& object) {
    object.print(); // Laufzeitentscheidung
}
\end{lstlisting}

\subsection{Statischer und dynamischer Typ}
\begin{description}
  \item[Statischer Typ] Der Typ, der zur Compilezeit bekannt ist, z.\,B. der Typ eines Zeigers oder einer Referenz.
  \item[Dynamischer Typ] Der tatsächliche Objekttyp, auf den ein Zeiger oder eine Referenz zur Laufzeit zeigt.
\end{description}

\section{Grundstrukturen, Sichtbarkeiten und UML}

\subsection{\code{struct} vs. \code{class}}
\begin{center}
\begin{tabular}{lll}
\toprule
Konstrukt & Standardsichtbarkeit & Typischer Einsatz \\
\midrule
\code{struct} & \code{public} & einfache Datenstrukturen \\
\code{class} & \code{private} & gekapselte Objekte mit Methoden \\
\bottomrule
\end{tabular}
\end{center}

\subsection{UML-Sichtbarkeiten}
\begin{center}
\begin{tabular}{ll}
\toprule
Symbol & Bedeutung \\
\midrule
\code{+} & public \\
\code{-} & private \\
\code{\#} & protected \\
\code{\textasciitilde} & package \\
\bottomrule
\end{tabular}
\end{center}

\subsection{UML-Attributzusätze}
\begin{description}
  \item[\code{/Attribut}] Abgeleitetes Attribut; keine direkte Initialisierung nötig, da es berechnet wird.
  \item[\code{readOnly}] Wert darf nicht verändert werden.
  \item[\code{unique / nonunique}] Werte sind eindeutig bzw. nicht eindeutig.
  \item[\code{ordered / unordered}] Werte sind geordnet bzw. ungeordnet.
  \item[\code{composite}] Teil einer Komposition.
  \item[\code{bag}] ungeordnete Sammlung mit möglichen Duplikaten.
  \item[\code{seq / sequence}] geordnete Sammlung.
  \item[\code{union}] Vereinigung mehrerer Mengen.
  \item[\code{query}] lesende Operation ohne Zustandsänderung.
\end{description}

\subsection{Parameterrichtungen in UML}
\begin{center}
\begin{tabular}{ll}
\toprule
Richtung & Bedeutung \\
\midrule
\code{in} & Parameter wird gelesen \\
\code{out} & Parameter wird geschrieben \\
\code{inout} & Parameter wird gelesen und geschrieben \\
\code{return} & Parameter ist Rückgabewert \\
\bottomrule
\end{tabular}
\end{center}

\section{Allgemeine Syntax}

\subsection{Scope-Resolution-Operator}
Der Scope-Resolution-Operator \code{::} wird verwendet, um auf Elemente in Klassen, Namespaces oder globalen Bereichen zuzugreifen.

\begin{lstlisting}[caption={Zugriff auf einen Namen im Namespace std}]
std::cout << "Hallo Welt";
\end{lstlisting}

\section{SOLID-Prinzipien}
\begin{description}
  \item[S -- Single Responsibility Principle] Eine Klasse sollte genau eine klar abgegrenzte Verantwortung haben.
  \item[O -- Open/Closed Principle] Softwareeinheiten sollten offen für Erweiterung, aber geschlossen für Änderung sein.
  \item[L -- Liskov Substitution Principle] Objekte einer Basisklasse müssen durch Objekte abgeleiteter Klassen ersetzbar sein.
  \item[I -- Interface Segregation Principle] Schnittstellen sollten klein und spezifisch sein.
  \item[D -- Dependency Inversion Principle] Abhängigkeiten sollten von Abstraktionen statt von konkreten Implementierungen abhängen.
\end{description}

\section{Naming Conventions}

\begin{center}
\begin{tabular}{lll}
\toprule
Element & Mögliche Konventionen & Gewählte Konvention \\
\midrule
Variablen & \code{camelCase}, \code{lower\_snake\_case} & \code{lower\_snake\_case} \\
Konstanten & \code{UPPER\_SNAKE\_CASE} & \code{UPPER\_SNAKE\_CASE} \\
Funktionen & \code{camelCase}, \code{lower\_snake\_case} & \code{camelCase} \\
Klassen / Structs / Enums & \code{PascalCase} & \code{PascalCase} \\
\bottomrule
\end{tabular}
\end{center}

\section{Namespaces}

Namespaces werden genutzt, um Code logisch zu strukturieren und Namenskonflikte zu vermeiden.

\begin{lstlisting}[caption={Namespace-Struktur}]
namespace project {
namespace modul {
    void init();
    void disable();
}
}
\end{lstlisting}

\begin{lstlisting}[caption={Aufruf einer Funktion im Namespace}]
project::modul::init();
\end{lstlisting}

\section{Variablen und Konstanten}

\subsection{\code{auto}}
Mit \code{auto} leitet der Compiler den Typ automatisch aus dem Ausdruck rechts der Zuweisung ab.

\begin{lstlisting}[caption={Typableitung mit auto}]
auto number = 42;      // int
auto value = 3.14;     // double
auto text = "Hallo";   // const char*
\end{lstlisting}

\subsection{\code{std::optional}}
\code{std::optional} ist ein Container aus der C++-Standardbibliothek für Werte, die vorhanden sein können, aber nicht müssen.

\begin{lstlisting}[caption={Optionaler Rückgabewert}]
#include <optional>

std::optional<int> find_value(bool found) {
    if (found) {
        return 42;
    }
    return std::nullopt;
}
\end{lstlisting}

\subsection{\code{const}}
\code{const} kennzeichnet Werte, Parameter, Referenzen oder Methoden als unveränderlich.

\begin{lstlisting}[caption={Konstante Referenz als Funktionsparameter}]
void print(const std::string& text) {
    // text darf hier nicht verändert werden
}
\end{lstlisting}

\subsection{\code{constexpr}}
\code{constexpr} kennzeichnet Werte oder Funktionen, die zur Compilezeit berechenbar sein müssen.

\begin{lstlisting}[caption={Compilezeit-Konstante}]
constexpr int square(int x) {
    return x * x;
}

constexpr int result = square(5);
\end{lstlisting}

\subsection{\code{static}}
\code{static} kann je nach Kontext bedeuten, dass eine Variable oder Funktion an eine Klasse, Datei oder Lebensdauer gebunden ist. Bei Klassenmembern bedeutet es: Der Member gehört zur Klasse und nicht zu einem einzelnen Objekt.

\section{Fließkommazahlen vergleichen}

Fließkommazahlen werden binär approximiert. Dadurch entstehen Rundungsfehler, sodass direkte Gleichheitsvergleiche unzuverlässig sein können. Stattdessen vergleicht man mit einer Toleranz.

\begin{lstlisting}[caption={Vergleich mit Epsilon}]
#include <cmath>

const double EPSILON = 1e-9;

bool almost_equal(double a, double b) {
    return std::fabs(a - b) < EPSILON;
}
\end{lstlisting}

\important{Merke:} \code{double} hat ungefähr 15 signifikante Dezimalstellen. Unterschiede kleiner als etwa \code{1e-15} können maschinelles Rauschen sein.

\section{Strings}

\subsection{Einbindung und Initialisierung}
\begin{lstlisting}[caption={String initialisieren}]
#include <string>

std::string name = "Content";
\end{lstlisting}

\subsection{Operatoren}
\begin{center}
\begin{tabular}{lll}
\toprule
Operation & Syntax & Bedeutung \\
\midrule
Zuweisung & \code{s1 = s2} & Inhalt zuweisen \\
Verkettung & \code{s1 += s2} & an bestehenden String anhängen \\
Verkettung & \code{s = s1 + s2} & neuen String bilden \\
Vergleich & \code{s1 == s2} & Gleichheit prüfen \\
\bottomrule
\end{tabular}
\end{center}

\subsection{Wichtige Methoden}
\begin{lstlisting}[caption={String-Methoden}]
std::string text = "Hallo Welt";

std::size_t length = text.length();
std::size_t position = text.find("Welt");
std::string part = text.substr(0, 5);
\end{lstlisting}

\begin{description}
  \item[\code{find(...)}] sucht von links.
  \item[\code{rfind(...)}] sucht von rechts.
  \item[\code{find\_first\_of(...)}] sucht das erste Zeichen aus einer Menge.
  \item[\code{find\_last\_of(...)}] sucht das letzte Zeichen aus einer Menge.
  \item[\code{find\_first\_not\_of(...)}] sucht das erste Zeichen, das nicht in der Menge ist.
  \item[\code{substr(start, length)}] extrahiert einen Teilstring.
\end{description}

\section{Casts}

\important{Grundregel:} Casts sollten bewusst und sparsam eingesetzt werden. Häufig sind sie ein Hinweis auf schlechtes Design oder fehlende Abstraktion.

\begin{center}
\begin{tabular}{ll}
\toprule
Cast & Zweck \\
\midrule
\code{static\_cast} & normale, explizite Umwandlungen, z.\,B. \code{int} zu \code{float} \\
\code{dynamic\_cast} & sichere Casts in Vererbungshierarchien \\
\code{const\_cast} & Entfernen oder Hinzufügen von \code{const} \\
\code{reinterpret\_cast} & sehr unsichere Low-Level-Umdeutung; möglichst vermeiden \\
\bottomrule
\end{tabular}
\end{center}

\section{Pointer und Referenzen}

\subsection{Referenzen}
Eine Referenz ist ein Alias für eine bestehende Variable.

\begin{lstlisting}[caption={Referenz}]
int a = 10;
int& ref = a;
\end{lstlisting}

\subsection{Pointer}
Ein Pointer speichert die Adresse einer Variable. Mit \code{*} wird dereferenziert.

\begin{lstlisting}[caption={Pointer und Dereferenzierung}]
int a = 10;
int* pointer = &a;

std::cout << pointer;   // Adresse von a
std::cout << *pointer;  // Wert von a
\end{lstlisting}

\subsection{Übergabe an Funktionen}

\subsubsection{Call-by-Value}
Die Funktion erhält eine Kopie. Änderungen wirken sich nicht auf das Original aus.

\begin{lstlisting}[caption={Call-by-Value}]
void change(int x) {
    x = 99;
}
\end{lstlisting}

\subsubsection{Call-by-Pointer}
Die Adresse wird übergeben. Das Original kann verändert werden. Achtung: Ein Pointer kann \code{nullptr} sein.

\begin{lstlisting}[caption={Call-by-Pointer}]
void change(int* p) {
    if (p != nullptr) {
        *p = 99;
    }
}
\end{lstlisting}

\subsubsection{Call-by-Reference}
Die Variable wird als Alias übergeben. Es entsteht keine Kopie.

\begin{lstlisting}[caption={Call-by-Reference}]
void change(int& ref) {
    ref = 99;
}
\end{lstlisting}

\subsubsection{Konstante Referenz}
Best Practice für lesenden Zugriff auf größere Objekte.

\begin{lstlisting}[caption={Konstante Referenz}]
void print(const std::string& text) {
    std::cout << text;
}
\end{lstlisting}

\section{Dynamische Speicherverwaltung}

\subsection{\code{new}}
\code{new} reserviert Speicher, erzeugt ein Objekt und ruft gegebenenfalls den Konstruktor auf.

\begin{lstlisting}[caption={Speicher mit new reservieren}]
int* p_int = new int;
int* int_initialisiert = new int(42);
int* int_array = new int[10];

MyClass* p_object = new MyClass();
\end{lstlisting}

\subsection{\code{delete}}
\code{delete} gibt Speicher frei, der mit \code{new} reserviert wurde. Bei Objekten wird zusätzlich der Destruktor aufgerufen.

\begin{lstlisting}[caption={Speicher freigeben}]
delete p_int;
delete p_object;
delete[] int_array;
\end{lstlisting}

\important{Hinweis:} In modernem C++ sollten häufig Smart Pointer wie \code{std::unique\_ptr} oder \code{std::shared\_ptr} statt manuellem \code{new}/\code{delete} verwendet werden.

\subsection{Smart Pointer}
Smart Pointer verwalten dynamisch erzeugte Objekte automatisch. Das zentrale Prinzip dahinter heißt \important{RAII} (Resource Acquisition Is Initialization): Eine Ressource wird an ein Objekt gebunden und automatisch freigegeben, sobald dieses Objekt seinen Gültigkeitsbereich verlässt.

Dafür muss der Header \code{<memory>} eingebunden werden.

\begin{lstlisting}[caption={Header für Smart Pointer}]
#include <memory>
\end{lstlisting}

\subsubsection{Warum Smart Pointer?}
Smart Pointer vermeiden typische Fehler der manuellen Speicherverwaltung:
\begin{itemize}
  \item vergessene \code{delete}-Aufrufe,
  \item doppelte Freigabe desselben Speichers,
  \item Speicherlecks bei Exceptions,
  \item unklare Besitzverhältnisse von Objekten.
\end{itemize}

\subsubsection{\code{std::unique\_ptr}}
\code{std::unique\_ptr} besitzt ein Objekt exklusiv. Es gibt genau einen Besitzer. Der Pointer kann nicht kopiert, aber verschoben werden.

\begin{lstlisting}[caption={Exklusiver Besitz mit std::unique\_ptr}]
#include <memory>

class Motor {
public:
    void start() {}
};

int main() {
    auto motor = std::make_unique<Motor>();
    motor->start();

    // Kein delete notwendig.
    // Der Motor wird automatisch freigegeben,
    // sobald motor den Scope verlässt.
}
\end{lstlisting}

\important{Best Practice:} Verwende \code{std::make\_unique<T>(...)} statt \code{new}, weil der Code sicherer und lesbarer ist.

\subsubsection{\code{std::shared\_ptr}}
\code{std::shared\_ptr} erlaubt gemeinsamen Besitz. Mehrere Smart Pointer können dasselbe Objekt besitzen. Das Objekt wird erst gelöscht, wenn der letzte \code{std::shared\_ptr} zerstört wurde.

\begin{lstlisting}[caption={Gemeinsamer Besitz mit std::shared\_ptr}]
#include <memory>

class Sensor {
public:
    void read() const {}
};

int main() {
    auto sensor_1 = std::make_shared<Sensor>();
    auto sensor_2 = sensor_1; // Referenzzähler wird erhöht

    sensor_1->read();
    sensor_2->read();

    // Das Objekt lebt, solange mindestens ein shared_ptr darauf zeigt.
}
\end{lstlisting}

\important{Best Practice:} Verwende \code{std::shared\_ptr} nur, wenn wirklich mehrere Besitzer notwendig sind. Sonst ist \code{std::unique\_ptr} meistens die bessere Wahl.

\subsubsection{\code{std::weak\_ptr}}
\code{std::weak\_ptr} beobachtet ein Objekt, das von \code{std::shared\_ptr} verwaltet wird, besitzt es aber nicht. Dadurch wird der Referenzzähler nicht erhöht. Das ist besonders wichtig, um zyklische Referenzen zu vermeiden.

\begin{lstlisting}[caption={Beobachtender Zugriff mit std::weak\_ptr}]
#include <memory>
#include <iostream>

int main() {
    auto shared = std::make_shared<int>(42);
    std::weak_ptr<int> weak = shared;

    if (auto locked = weak.lock()) {
        std::cout << *locked << '\n';
    }
}
\end{lstlisting}

\subsubsection{Vergleich der Smart Pointer}
\begin{center}
\begin{tabular}{p{3.2cm}p{5.4cm}p{5.2cm}}
\toprule
Smart Pointer & Besitzmodell & Typischer Einsatz \\
\midrule
\code{std::unique\_ptr} & exklusiver Besitz; nicht kopierbar, nur verschiebbar & Standardwahl für dynamische Objekte mit eindeutigem Besitzer \\
\code{std::shared\_ptr} & gemeinsamer Besitz mit Referenzzählung & mehrere unabhängige Besitzer benötigen dasselbe Objekt \\
\code{std::weak\_ptr} & kein Besitz; beobachtet ein \code{std::shared\_ptr}-Objekt & Vermeidung zyklischer Referenzen und optionaler Zugriff \\
\bottomrule
\end{tabular}
\end{center}

\subsubsection{Merksätze}
\begin{itemize}
  \item Verwende \code{std::unique\_ptr}, wenn genau ein Objekt für die Lebensdauer verantwortlich ist.
  \item Verwende \code{std::shared\_ptr}, wenn sich mehrere Stellen den Besitz wirklich teilen müssen.
  \item Verwende \code{std::weak\_ptr}, wenn du ein Objekt nur beobachten willst.
  \item Verwende \code{std::make\_unique} und \code{std::make\_shared} statt direktem \code{new}.
  \item Manuelles \code{delete} sollte in modernem C++ möglichst selten notwendig sein.
\end{itemize}

\section{Klassen}

\subsection{Sichtbarkeiten}
\begin{description}
  \item[\code{private}] Methoden und Attribute sind nur innerhalb der Klasse sichtbar. Ausnahme: \code{friend}.
  \item[\code{protected}] Sichtbar in der Klasse selbst und in abgeleiteten Klassen.
  \item[\code{public}] Öffentlich sichtbar.
\end{description}

\subsection{Konstruktor und Destruktor}
Konstruktoren initialisieren Objekte. Destruktoren räumen Ressourcen beim Zerstören eines Objekts auf.

\begin{lstlisting}[caption={Konstruktor und Destruktor}]
class Klassenname {
public:
    Klassenname();   // Konstruktor
    ~Klassenname();  // Destruktor
};
\end{lstlisting}

\subsection{Attribute initialisieren}
\subsubsection{Problematische Zuweisung}
\begin{lstlisting}[caption={Unklare Initialisierung}]
Auto::Auto(std::string ps, int seats) {
    ps = ps;
    seats = seats;
}
\end{lstlisting}

\subsubsection{Zuweisung mit \code{this}}
\begin{lstlisting}[caption={Zuweisung mit this}]
Auto::Auto(std::string ps, int seats) {
    this->ps = ps;
    this->seats = seats;
}
\end{lstlisting}

\subsubsection{Bessere Lösung: Initialisierungsliste}
\begin{lstlisting}[caption={Initialisierungsliste}]
Auto::Auto(std::string ps, int seats)
    : ps{ps}, seats{seats} {
}
\end{lstlisting}

\subsection{Konstante Memberfunktionen}
Konstante Memberfunktionen dürfen keine Attribute des Objekts verändern. Sie eignen sich für Query-Funktionen.

\begin{lstlisting}[caption={Konstante Memberfunktion}]
class Auto {
public:
    int get_ps() const {
        return ps;
    }

private:
    int ps;
};
\end{lstlisting}

\section{Statische Klassenmember}

\subsection{Statische Variablen}
Statische Klassenmember sind an die Klasse gebunden und nicht an ein einzelnes Objekt. Sie existieren nur einmal.

\begin{lstlisting}[caption={Statische Klassenvariable}]
class Auto {
public:
    static int anzahl_autos;

    Auto() {
        ++anzahl_autos;
    }
};

int Auto::anzahl_autos = 0;
\end{lstlisting}

\subsection{Statische Methoden}
Statische Methoden funktionieren ohne Objekt. Sie besitzen keinen \code{this}-Pointer und können direkt nur auf statische Member zugreifen.

\begin{lstlisting}[caption={Statische Methode}]
class Rechner {
public:
    static int add(int a, int b) {
        return a + b;
    }
};

int ergebnis = Rechner::add(5, 5);
\end{lstlisting}

\section{Vererbung}

\subsection{Basisklasse und abgeleitete Klasse}
\begin{lstlisting}[caption={Einfache Vererbung}]
class Base {
public:
    Base() = default;
};

class Child : public Base {
public:
    Child() = default;
};
\end{lstlisting}

\subsection{Aufruf des Basisklassen-Konstruktors}
Wenn die Basisklasse keinen Default-Konstruktor besitzt, muss die abgeleitete Klasse den passenden Konstruktor explizit aufrufen.

\begin{lstlisting}[caption={Basisklassen-Konstruktor aufrufen}]
class Base {
public:
    explicit Base(int value);
};

class Child : public Base {
public:
    Child(int x) : Base{x} {
    }
};
\end{lstlisting}

\subsection{Virtueller Destruktor}
Bei Polymorphie sollte die Basisklasse einen virtuellen Destruktor besitzen.

\begin{lstlisting}[caption={Virtueller Destruktor}]
class Base {
public:
    virtual ~Base() = default;
};

class Child : public Base {
public:
    ~Child() override = default;
};
\end{lstlisting}

\section{Klassenbeziehungen}

\subsection{Assoziation}
Eine Assoziation ist eine lockere Beziehung zwischen zwei Klassen.

\begin{lstlisting}[caption={Assoziation}]
class Fahrer;

class Auto {
public:
    Fahrer* fahrer;
};
\end{lstlisting}

\subsection{Aggregation}
Eine Aggregation ist ebenfalls eine lockere Beziehung. Das enthaltene Objekt kann unabhängig vom Besitzer existieren.

\subsection{Komposition}
Eine Komposition ist eine starke, exklusive Beziehung. Das Teilobjekt gehört fest zum Ganzen.

\begin{lstlisting}[caption={Komposition}]
class Rad {};

class Auto {
public:
    Rad raeder[4];
};
\end{lstlisting}

\section{STL-Container}

\subsection{Wichtige Container}
\begin{itemize}
  \item \code{std::vector}
  \item \code{std::list}
  \item \code{std::deque}
  \item \code{std::queue}
  \item \code{std::priority\_queue}
  \item \code{std::map}
\end{itemize}

\subsection{Gemeinsame Funktionen}
\begin{center}
\begin{tabular}{ll}
\toprule
Funktion & Bedeutung \\
\midrule
\code{x.size()} & Anzahl der Elemente \\
\code{x.empty()} & \code{true}, wenn Container leer ist \\
\code{x.insert(position, wert)} & Wert einfügen \\
\code{x.erase(position)} & Wert löschen \\
\bottomrule
\end{tabular}
\end{center}

\section{Iteratoren}

\subsection{Iterator-Kategorien}
\begin{center}
\begin{tabular}{lll}
\toprule
Kategorie & Eigenschaften & Beispiele \\
\midrule
Input & liest Daten einmalig, nur vorwärts & Input Streams \\
Output & schreibt Daten, nur vorwärts & Output Streams \\
Forward & liest/schreibt mehrfach, nur vorwärts & \code{forward\_list}, \code{unordered\_map} \\
Bidirectional & vorwärts und rückwärts & \code{std::list}, \code{std::map} \\
Random Access & beliebiger Zugriff, pointerähnlich & \code{std::vector} \\
\bottomrule
\end{tabular}
\end{center}

\subsection{Iterator-Funktionen}
\begin{description}
  \item[\code{x.begin()}] Iterator auf das erste Element.
  \item[\code{x.end()}] Iterator hinter das letzte Element.
\end{description}

\begin{lstlisting}[caption={Iterator-Beispiel}]
std::vector<int> values{1, 2, 3};

for (std::vector<int>::iterator it = values.begin(); it != values.end(); ++it) {
    std::cout << *it << " ";
}
\end{lstlisting}

\section{STL-Algorithmen}

\subsection{Suchalgorithmen}
\begin{description}
  \item[\code{std::find}] findet das erste passende Element.
  \item[\code{std::search}] sucht ein Muster.
\end{description}

\subsection{Sortieralgorithmen}
\begin{description}
  \item[\code{std::sort}] sortiert Elemente aufsteigend.
  \item[\code{std::reverse}] kehrt die Reihenfolge um.
  \item[\code{std::shuffle}] mischt Elemente zufällig.
\end{description}

\subsection{Numerische Algorithmen}
\begin{description}
  \item[\code{std::accumulate}] berechnet z.\,B. eine Summe.
  \item[\code{std::count}] zählt Vorkommen eines bestimmten Wertes.
\end{description}

\subsection{Modifizierende Algorithmen}
\begin{description}
  \item[\code{std::copy}] kopiert Elemente.
  \item[\code{std::transform}] wendet eine Funktion auf Elemente an.
  \item[\code{std::replace}] ersetzt Werte.
\end{description}

\begin{lstlisting}[caption={std::transform mit Lambda}]
std::transform(values.begin(), values.end(),
               values.begin(),
               [](int x) {
                   return x * 2;
               });
\end{lstlisting}

\section{Range-based for-Schleifen}

\subsection{Kopie}
\begin{lstlisting}[caption={for-each mit Kopie}]
for (int x : container) {
    // x ist eine Kopie
}
\end{lstlisting}
Änderungen an \code{x} wirken sich nicht auf den Container aus.

\subsection{Referenz}
\begin{lstlisting}[caption={for-each mit Referenz}]
for (int& x : container) {
    // x ist der Originalwert
}
\end{lstlisting}
Änderungen an \code{x} wirken sich auf den Container aus.

\subsection{Konstante Referenz}
\begin{lstlisting}[caption={Best Practice: konstante Referenz}]
for (const int& x : container) {
    // effizient und sicher
}
\end{lstlisting}

\subsection{Structured Bindings}
\begin{lstlisting}[caption={Structured Binding bei std::map}]
std::map<double, int> m;

for (const auto& [key, value] : m) {
    std::cout << key << ": " << value << '\n';
}
\end{lstlisting}

\important{Hinweis:} Die Syntax \code{[key, value]} ist ein Structured Binding und kein Array.

\section{Templates}

Templates ermöglichen generische Funktionen und Klassen. Sie werden zur Compilezeit instanziiert. Der Compiler erzeugt für die verwendeten Typen passende konkrete Versionen.

\begin{lstlisting}[caption={Generische Funktion mit Template}]
template <typename T, int SIZE>
void print(T (&array)[SIZE]) {
    for (int i = 0; i < SIZE; ++i) {
        std::cout << array[i] << '\n';
    }
}
\end{lstlisting}

\subsection{Eigenschaften von Templates}
\begin{itemize}
  \item Sie werden zur Compilezeit aufgelöst.
  \item Sie können für verschiedene Typen verwendet werden.
  \item Sie werden in jeder Übersetzungseinheit separat instanziiert, wenn sie dort verwendet werden.
  \item Spezialisierungen können für Sonderfälle definiert werden.
\end{itemize}

\section{Exception Handling}

Exception Handling dient der Fehlerbehandlung mit Ausnahmen. Fehler können gemeldet, abgefangen und behandelt werden.

\subsection{Grundbegriffe}
\begin{description}
  \item[\code{throw}] meldet einen Fehler.
  \item[\code{try}] markiert einen Codeblock, in dem Fehler auftreten können.
  \item[\code{catch}] fängt und behandelt Fehler.
\end{description}

\subsection{Wichtige Standard-Exceptions}
\begin{center}
\begin{tabular}{ll}
\toprule
Exception & Bedeutung \\
\midrule
\code{std::invalid\_argument} & ungültiges Argument \\
\code{std::out\_of\_range} & Index außerhalb des erlaubten Bereichs \\
\code{std::runtime\_error} & allgemeiner Laufzeitfehler \\
\code{std::bad\_alloc} & Speicherallokation fehlgeschlagen \\
\bottomrule
\end{tabular}
\end{center}

Alle genannten Exceptions erben von \code{std::exception}. Mit \code{error.what()} kann eine Fehlermeldung ausgegeben werden.

\subsection{Catch-Reihenfolge}
Spezifische Exceptions sollten vor allgemeinen Exceptions gefangen werden.

\begin{lstlisting}[caption={try-catch-Struktur}]
try {
    // riskanter Code
}
catch (const std::invalid_argument& error) {
    std::cout << "Ungueltiges Argument: " << error.what() << '\n';
}
catch (const std::out_of_range& error) {
    std::cout << "Bereichsfehler: " << error.what() << '\n';
}
catch (const std::exception& error) {
    std::cout << "Standard-Exception: " << error.what() << '\n';
}
catch (...) {
    std::cout << "Unbekannter Fehler\n";
}
\end{lstlisting}

\subsection{Best Practice}
\begin{lstlisting}[caption={Allgemeiner catch-Block}]
catch (const std::exception& e) {
    std::cout << e.what() << '\n';
}
\end{lstlisting}

\subsection{\code{noexcept}}
Eine Funktion, die keine Exceptions werfen darf, kann mit \code{noexcept} markiert werden.

\begin{lstlisting}[caption={noexcept-Funktion}]
void foo() noexcept {
}
\end{lstlisting}

\end{document}
